\scanpage{189}
\mspage{497}
% position: top continuation
\texttt{ou même dans le monde chinois (ou plus généralement le monde ``oriental'') d'au-
jourd'hui. Aucun signe, au niveau de la vie de tous les jours, ne pourrait me le}
faire supposer, ni à travers mes amis et amies orientaux, ni à travers les échos
qui ont pu me parvenir de la tradition et de la vie d' aujourd'hui dans la
Chine ou d'autres pays d'extrême-Orient - bien au contraire. Il me semble plutôt
qu'une fine perception du dynamisme yin-yang s'est cantonnée à peu près exclu-
sivement dans la \mbox{p r a t i q u e} de \mbox{c e r t a i n s\quad a r t s} - comme
la calligraphie, la poésie, l'art culinaire et, bien entendu, l'art médical (*).

% position: middle paragraph
C'est le dernier surtout, sous le nom de ``médecine chinoise'' et par le
biais de certains succés spectaculaires de l'acupuncture, qui a fini au cours \unsure{des}
vingt ans passés, par acquérir droit de cité chez nous, et à être investi de
prestige. Nombreux pourtant sont encore ceux qui ignorent qu'en médecine chi-
noise, l'alpha et l'oméga de l'appréhension du corps, de la circulation d'éner-
gie dans le corps et des perturbations de celle-ci (qui constituent les états
morbides que nous appelons ``maladies''), se trouve justement dans une dialecti-
que très fine du yin et du yang. Le fait que cette dialectique ``marche'', puis-
que la ``médecine chinoise'' basée sur elle est efficace (y compris dans de nom-
breux cas qui échappent aux moyens de la panoplie occidentale), peut être con-
sidéré comme une sorte de ``preuve'' de la réalité des ``principes'' ou ``aspects''
ou ``modes'' (d'appréhension, ou d'existence) yin et yang - que ce ne sont pas
de pures spéculations sorties des chapeaux de certains philosophes et autres
poètes (pour ne pas dire fumistes).

% position: lower paragraph
On peut se demander, il est vrai, quel est le sens de telles preuves, et
même de toute ``preuve'' quelle qu'elle soit de la validité de telle ou telle
vision du monde. A supposer même que la preuve ait convaincu (c'est-à-dire,
que l'intéressé ait bien voulu se laisser convaincre), et même et par dessus
le marché\unsure{\inserted{,}} que la vision en question soit profonde et par là\unsure{\inserted{,}} bienfaisante -
la meilleure preuve du monde est impuissante pourtant à \mbox{c o m m u n i q u e r}
% position: footnote
\par\noindent\rule{0.24\linewidth}{0.3pt}\par
(*) (21 octobre) J'ai oublié dans le nombre l'\mbox{a r t\quad d i v i n a t o i r e},
dans le \mbox{Y i\quad K i n g} ou ``livre des changes'', qui jouit aujourd'hui d'une\unsure{.} grande
popularité dans certains milieux tant en Europe qu'en Amérique. Les \unsure{$\cdot$}64 ``hexagram\unsure{-}
mes'' qui constituent les ``mots'' de base du langage divinatoire du Yi King, ne \unsure{son|sont}
autres que les $2^{\inserted{6}}$ combinaisons possibles de suites de six ``signes'' yin et yang,
depuis le yin pur (six répétitions du yin ) au yang pur (six répétitions du \unsure{yang|yang)}
Il semble y avoir là une sorte d'alchimie d'une grande finesse des combinaisons
du yin et du yang, qui \inserted{(}paraît-il\inserted{)} avait fasciné Jung. L'intérêt de cette alchimie
\inserted{(}en tant que ``collection d'archétypes'' notamment) me paraît a priori indépendant
de son usage en art divinatoire, et du crédit qu'on est disposé à accorder à
un tel usage.
\margin{188}
